\chapter{Prefacio}

Este libro nace de una ambición sencilla: un único curso de física,
coherente de principio a fin, que abarque desde la educación infantil
hasta el final del grado y que pueda leerse en cualquier parte del
mundo.

El estilo es conciso y riguroso: un curso construido a partir de
definiciones, ejemplos, proposiciones y teoremas, con demostraciones
siempre que resulten accesibles al nivel considerado, seguido de
ejercicios cuidadosamente elegidos. Los resultados que se enuncian sin
demostración se señalan explícitamente como admitidos. Las soluciones
completas de todos los ejercicios se recogen al final del libro.

Cada curso escolar forma una parte del libro, y cada parte se divide en
capítulos. Cuanto más temprano es el curso, más joven es el lector al que
se dirige: más ilustraciones, más pasos detallados y ejercicios más
suaves.

\section*{Cómo leer este libro}

Los enunciados van numerados dentro de cada capítulo y se distinguen por
colores: las definiciones en azul, los teoremas en rojo, las
proposiciones y sus parientes en naranja, y los métodos --- recetas
breves para las tareas habituales --- en verde. Los ejercicios están
graduados desde ($\star$), de aplicación directa, hasta
($\star\star\star$), los problemas más difíciles.

\section*{Cómo contribuir}

Las fuentes de este libro están en un repositorio git público:
\begin{center}
\url{https://github.com/bvirrion/one-physics-book}
\end{center}
Toda contribución --- nuevos capítulos, correcciones, mejores
demostraciones, ejercicios adicionales --- es bienvenida; véase
\texttt{CONTRIBUTING.md} en la raíz del repositorio.
